\documentclass{article}

\usepackage{amsmath,amssymb}
\usepackage{xurl}
\usepackage[hidelinks]{hyperref}
\setlength{\emergencystretch}{2em}

% Shared commands for notes in docs/paper-gaps/.

\DeclareUnicodeCharacter{2080}{\ifmmode{}_{0}\else\textsubscript{0}\fi}
\DeclareUnicodeCharacter{2081}{\ifmmode{}_{1}\else\textsubscript{1}\fi}
\DeclareUnicodeCharacter{2082}{\ifmmode{}_{2}\else\textsubscript{2}\fi}
\DeclareUnicodeCharacter{2099}{\ifmmode{}_{n}\else\textsubscript{n}\fi}

\newcommand{\Lean}{\textsc{Lean}}
\ExplSyntaxOn
\NewDocumentCommand{\leanid}{m}{
  \ifmmode
    \textsf{\detokenize{#1}}
  \else
    \group_begin:
    \urlstyle{sf}
    \tl_set:Nn \l_tmpa_tl {#1}
    \tl_replace_all:Nnn \l_tmpa_tl {\_} {_}
    \exp_args:NV \path \l_tmpa_tl
    \group_end:
  \fi
}
\ExplSyntaxOff
\providecommand{\tr}{\operatorname{tr}}
\newcommand{\tnleanIssueBase}{https://github.com/LionSR/TNLean/issues/}
\newcommand{\tnleanPrBase}{https://github.com/LionSR/TNLean/pull/}
\newcommand{\tnleanDocBase}{https://sirui-lu.com/TNLean/docs/}
\newcommand{\ghissue}[1]{\href{\tnleanIssueBase#1}{GitHub issue \##1}}
\newcommand{\ghpr}[1]{\href{\tnleanPrBase#1}{PR \##1}}
\newcommand{\doclink}[1]{\href{\tnleanDocBase#1.html}{\path{#1}}}

% Dirac notation and the identity operator, as in blueprint/src/macros/common.tex.
% \providecommand keeps note-local definitions authoritative.
\usepackage{dsfont}
\providecommand{\ket}[1]{|#1\rangle}
\providecommand{\bra}[1]{\langle #1|}
\providecommand{\braket}[2]{\langle #1 | #2 \rangle}
\providecommand{\ketbra}[2]{|#1\rangle\!\langle #2|}
\providecommand{\Id}{\mathds{1}}
\providecommand{\one}{\mathds{1}}
\providecommand{\C}{\mathbb{C}}
\providecommand{\MN}[1]{M_{#1}(\C)}
\providecommand{\MD}{M_D(\C)}

% External referencing for paper sources.  The build helper writes sanitized
% auxiliary files under build/paper-aux/ and excludes duplicated source labels.
\usepackage{xr}

% Import a generated paper auxiliary file with a caller-supplied label prefix.
% The candidate paths support builds from docs/paper-gaps/, the repository root,
% and build/paper-aux/ itself.
\newcommand{\importpaperaux}[2]{%
  \IfFileExists{../../build/paper-aux/#2.aux}{%
    \externaldocument[#1]{../../build/paper-aux/#2}%
  }{%
    \IfFileExists{build/paper-aux/#2.aux}{%
      \externaldocument[#1]{build/paper-aux/#2}%
    }{%
      \IfFileExists{#2.aux}{%
        \externaldocument[#1]{#2}%
      }{}%
    }%
  }%
}

% General external reference macros
\newcommand{\paperref}[2]{\ref{#1-#2}}
\newcommand{\papereqref}[2]{(\ref{#1-#2})}

% CPSV16 source (1606.00608 / MPDO-22-12-17-2.tex) external document and shortcuts
\importpaperaux{cpsv16-}{1606.00608/MPDO-22-12-17-2-tnlean}
\newcommand{\cpsvref}[1]{\paperref{cpsv16}{#1}}
\newcommand{\cpsveqref}[1]{\papereqref{cpsv16}{#1}}

% Verdict marker: \gapnote{<kind>}{<status>}. Kind is the policy
% classification (clarification, local-correction, scope-restriction,
% unfaithful, false-source, open-gap); status is open, wip (elimination
% underway), resolved, or historical. Machine-read by the site index;
% prints a verdict line.
\newcommand{\gapnote}[2]{\par\noindent\textbf{Verdict:} #1 (#2).\par}

\providecommand{\ghissue}[1]{GitHub issue~\##1}
\providecommand{\leanid}[1]{\path{#1}}

\title{Normal PEPS Fundamental Theorem: Section 3 Route}
\author{}
\date{2026-05-06}

\begin{document}
\maketitle

\section{The source argument}

The normal-PEPS part of \cite{Molnar2018NormalPEPS} begins after the injective
PEPS theorem. Its purpose is not to replace the injective proof, but to make
the injective proof applicable after a controlled blocking. A normal PEPS is
one for which sufficiently large regions can be blocked into injective tensors.
Once such injective regions are available, the proof again applies the
three-site injective-chain argument around a chosen edge.

The source proof has three layers. First, Lemma \path{injective_union} proves
that the union of two injective regions is injective. The proof blocks the
network into the four regions $A\setminus B$, $A\cap B$,
$B\setminus A$, and $(A\cup B)^c$. If a boundary tensor kills the
contraction over $A\cup B$, injectivity of $A$ first removes the
$A\setminus B$ and $A\cap B$ part, and injectivity of $B$ then removes
the remaining $A\cap B$ and $B\setminus A$ part.
\footnote{The source is
    \path{Papers/1804.04964/paper_normal.tex:1324--1400}.}

Second, for the translationally invariant square-lattice theorem, the paper
assumes that every $2\times 3$ and $3\times 2$ region is injective. By
repeatedly applying the union lemma, the regions denoted $R$, $S$, and $T$
in the paper are injective. The finite-lattice block denoted $A_3$ in the proof
of Theorem~\path{3} has the same complementary shape as $T$ but lives in the
full finite lattice around the chosen edge; the present coordinate definition
records only the displayed local-window complement. The regions $R$ and $S$
differ by one site, and their complements are injective when the system is
large enough. These are the regions used in the last comparison step.
\footnote{The regions $R$, $S$, and $T$ are drawn at
    \path{Papers/1804.04964/paper_normal.tex:1407--1443}.}

Third, the proof of Theorem \path{3} blocks a neighborhood of each edge into
three injective pieces. The red region is $A_1$, the blue region is $A_2$,
and the complement is $A_3$. This produces the same three-site injective MPS
configuration used in the injective PEPS proof. Lemma \path{inj_isomorph}
therefore gives a gauge on every edge. Translation invariance reduces these
edge gauges to one horizontal matrix $X$ and one vertical matrix $Y$. After
absorbing them into the second tensor, the final comparison of the $R$ and
$S$ regions gives $A\propto \widetilde B$.
\footnote{The proof of Theorem \path{3} is
    \path{Papers/1804.04964/paper_normal.tex:1449--1572}.}

The general theorem \path{normal} removes the square-lattice specialization.
It assumes precisely the blockings needed above: every edge can be blocked into
a three-site injective chain, and every site admits two injective regions with
injective complements that differ exactly in that site.
\footnote{The general theorem is
    \path{Papers/1804.04964/paper_normal.tex:1574--1585}.}

\section{Current formal representation}

The injective PEPS objects used by the normal theorem are represented in
\path{TNLean/PEPS/Defs.lean}, \path{TNLean/PEPS/VirtualInsertion.lean}, and
\path{TNLean/PEPS/Blocking.lean}. Chapter 13a of the blueprint has a
separate normal-PEPS section after the injective theorem. This separation is
important: the normal theorem needs extra blocking hypotheses before the
injective proof route can be invoked.

The present blueprint nodes are the following.

\begin{itemize}
    \item \path{lem:peps_injectiveRegion_union} states the union-of-injective
          regions lemma. It is tracked by \ghissue{1374}. It records the
          four-region decomposition and the two successive inverse applications
          used in the source proof.
    \item \path{lem:peps_regionInjectivityUnionClosure_finite} records the
          finite nonempty iteration of the binary union-closure assertion.
    \item \path{def:peps_normalSquareBlockingRegions} records the square-lattice
          regions $R$, $S$, and $T$. It is part of \ghissue{1375}. The
          coordinate rectangular hypotheses, union closure, and a supplied
          rectangular cover of the displayed $T$-region imply this
          abstract structure.
    \item \path{thm:peps_normalEdgeBlockingToInjectiveChain} states that the
          normal square-lattice hypotheses allow a blocking around every edge into
          three injective regions. It is part of \ghissue{1375}.
    \item \path{thm:peps_tiNormalGaugeAbsorption} records the translationally
          invariant gauge absorption with horizontal gauge $X$ and vertical gauge
          $Y$. It is part of \ghissue{1375}.
    \item \path{thm:fundamentalTheorem_normalSquarePEPS} records the square-lattice
          theorem corresponding to Theorem \path{3}. It is part of \ghissue{1375}.
    \item \path{thm:fundamentalTheorem_normalPEPS} records the general normal PEPS
          theorem following Theorem \path{3}. It is part of \ghissue{1375}.
\end{itemize}

The finite rectangular square lattice is represented by pairs of horizontal
and vertical coordinates. Its graph has horizontal and vertical
nearest-neighbor edges, named coordinate right/up edges, and orientation
predicates distinguishing the two possible directions. Every edge has exactly
one of these orientations. With the ordered endpoints convention for edges,
horizontal edges are coordinate right edges from their left endpoints, and
vertical edges are coordinate upward edges from their lower endpoints.

The four-region decomposition from Lemma \path{injective_union} is formalized
both by named regions and by an indexed family in the source-paper order:
$A\setminus B$, $A\cap B$, $B\setminus A$, and $(A\cup B)^c$. The indexed
family reconstructs $A$, $B$, $A\cup B$, $A^c$, $B^c$, and the full vertex set.
The same reconstructions are available as finite unions over the selected
indexed subfamilies $\{0,1\}$, $\{1,2\}$, $\{0,1,2\}$, $\{2,3\}$, and
$\{0,3\}$. The singleton subfamily $\{3\}$ gives the outside block and is the
complement of the selected inside subfamily $\{0,1,2\}$. Its selected
complementary pairs match the two places where the source proof applies
injectivity of $A$ and then injectivity of $B$. The
selected subfamilies $\{0,1\}$ and $\{2,3\}$ are disjoint, as are
$\{1,2\}$ and $\{0,3\}$; these are the disjointness forms needed for the two
inverse applications. The abstract union-closure hypothesis also has the
corresponding indexed inside-union form: injectivity of the selected subfamilies
$\{0,1\}$ and $\{1,2\}$ implies injectivity of the inside subfamily
$\{0,1,2\}$. Every vertex belongs to exactly one indexed piece. These facts
prepare the tensor-level boundary-contraction proof, but do not prove that
theorem.

The rectangular vocabulary separates coordinate products from source-paper
contiguous blocks. Coordinate rectangles are products of finite coordinate
sets, whereas contiguous coordinate rectangles are products of coordinate
intervals. The $2\times3$ and $3\times2$ coordinate-product predicates assert
only the two cardinalities; the contiguous-rectangle predicates express the
source-paper blocks. The square-lattice rectangular injectivity hypotheses use
these two contiguous predicates and map to the abstract rectangular-injectivity
bundle. Bounded contiguous rectangles of sizes $2\times3$ and $3\times2$
satisfy the corresponding source-rectangle predicates.

The displayed regions $R$ and $S$ are explicit unions of contiguous
rectangles. Region $R$ is the union of one contiguous $2\times3$ rectangle and
one contiguous $3\times2$ rectangle. Region $S$ is the union of two contiguous
$2\times3$ rectangles. Their coordinate definitions agree with the source
claim that they differ in one site: $R\subseteq S$, and $S\setminus R$ is the
lower-right site of the displayed $3\times3$ square. Assuming the
union-closure assertion supplied by the source lemma on unions of injective
regions, these decompositions imply injectivity of $R$ and $S$. The formal
development also contains the finite nonempty iteration of this union-closure
assertion, matching the repeated unions used after Lemma
\path{injective_union}.

The displayed local model for $T$ is the complement, inside the source
$5\times6$ window, of the two edge blocks in the source picture. These removed
blocks are one contiguous $2\times3$ rectangle and one contiguous $3\times2$
rectangle, hence they are injective under the coordinate rectangular
hypotheses. They are disjoint, and together with the local-window $T$ region
they reconstruct the $5\times6$ window. This local-window region is distinct
from the finite-lattice complementary block $A_3$ used later in the proof of
Theorem~\path{3}. The finite-lattice block is the complement of the two edge
blocks in the full square lattice; its elementary disjointness and covering
identities are formalized separately.

A target-parametric rectangular-cover criterion supplies a sufficient
injectivity route: if a square-lattice region has a nonempty cover by
contiguous $2\times3$ and $3\times2$ rectangles, then rectangular injectivity
and finite union closure imply injectivity of that region. The formal special
cases are the local-window $T$ region and the finite-lattice complementary
block. This criterion is only a sufficient local model, not the source
argument itself.

The origin-based local-window model for $T$ has no such rectangular cover
whenever the ambient lattice contains the corresponding $5\times6$ window.
The obstruction is a point-forcing argument. For the local-window target
$T_{\mathrm{loc}}$ at the origin, let $R^{a\times b}_{x,y}$ denote the
contiguous $a\times b$ rectangle with lower-left corner $(x,y)$. Then
\begin{align}
    p &:= (0,0) \in T_{\mathrm{loc}},                                    \\
    p \in R^{2\times3}_{x,y}
        &\Longrightarrow q_{23}:=(0,2)\in R^{2\times3}_{x,y},
        \qquad q_{23}\notin T_{\mathrm{loc}},                             \\
    p \in R^{3\times2}_{x,y}
        &\Longrightarrow q_{32}:=(2,1)\in R^{3\times2}_{x,y},
        \qquad q_{32}\notin T_{\mathrm{loc}}.
\end{align}
Thus any exact cover of $T_{\mathrm{loc}}$ by $2\times3$ and $3\times2$
rectangles would cover $p$ by one of its pieces and would therefore also cover
a point outside $T_{\mathrm{loc}}$. The same calculation applies to the
rotated local-window model for vertical edges in a $6\times5$ window, to the
origin-based horizontal edge-complement model in the $5\times7$ frame, and to
the origin-based vertical edge-complement model in the $7\times5$ frame. The
forced points are
\[
\begin{array}{c|c|c|c}
    \text{target} & p & q_{23} & q_{32} \\
    \hline
    T_{\mathrm{loc}} & (0,0) & (0,2) & (2,1) \\
    T_{\mathrm{loc}}^{\mathrm{rot}} & (0,0) & (1,2) & (2,1) \\
    A_{3,\mathrm{hor}} & (0,0) & (0,2) & (2,1) \\
    A_{3,\mathrm{ver}} & (0,0) & (1,2) & (2,1)
\end{array}
\]
The formal point-forcing lemma is
\leanid{TNLean.PEPS.not_squareLatticeRectangleCover_of_forced_points}. Its
four present applications are the corresponding no-cover theorems for the
local $T$ models and the edge-complement models. These no-cover results do not
contradict the paper. They show that the source's injectivity argument for $T$
and $A_3$ must be aligned with the finite PEPS geometry, rather than read
directly from the present exact-cover criterion.

The normalized horizontal edge frame has a collar identity: in the
$5\times7$ frame, $A_3$ is the local-window $T$ region together with two
top-collar contiguous $3\times2$ rectangles. Thus local $T$-injectivity implies
injectivity of $A_3$ by union closure and rectangular injectivity. The
vertical counterpart is a $7\times5$ frame in which the edge complement is the
rotated local $T$ region together with two right-collar contiguous
$2\times3$ rectangles. A rectangular cover of the rotated local $T$ region
therefore implies injectivity of the vertical complementary block.

The normalized horizontal and vertical edge blockings give one-edge data with
red, blue, and complementary regions. The distinguished edges have the
expected orientations; their endpoints lie in the red and blue regions; the
three regions are pairwise disjoint and cover the finite lattice. This one-edge
datum is the local form of the later every-edge blocking hypotheses. A family
of such one-edge data determines those hypotheses, and each datum supplies
the three injective regions used in the associated edge-blocked chain.

The translated horizontal and vertical pictures carry the same data at an
arbitrary coordinate origin. When the translated frame fits in the finite
lattice, the translated edge is identified with the corresponding coordinate
right or upward edge. A rectangular cover of the translated complementary
region supplies the one-edge blocking datum and hence the three injective
regions. The normalized pictures are the origin specializations of these
translated pictures, while the normalized and translated data remain in
separate Lean modules.

The conditional every-edge construction has two equivalent forms. A translated
edge-window assignment gives, for every edge, a horizontal or vertical window
and a rectangular cover of its complement. Such an assignment implies the
normal edge-blocking hypotheses. An oriented margin-and-cover datum records the
same input directly on each square-lattice edge: the edge orientation, the
coordinate inequalities placing the translated frame, and the complementary
cover. A family of these data also determines the normal edge-blocking
hypotheses, and the resulting hypotheses recover the supplied one-edge datum
at each edge.

For the present open rectangular coordinate graph, the margin predicates
describe only interior translated frames. A coordinate right edge must have
enough left and right margin; a coordinate upward edge must have enough bottom
and top margin. Hence neither the translated-window criterion nor the
oriented margin-and-cover criterion supplies data on all edges of the open
$7\times7$ graph. Explicit boundary obstructions are recorded on all four
sides. The remaining every-edge construction must therefore include the
boundary geometry and the source-faithful local and rotated $T$ covers.

Verdict: this note records an open gap with a scope restriction. The earlier
vocabulary gap for contiguous rectangular blocks is closed. The earlier
coordinate rectangular-injectivity structure is also present, and the
indexed finite-region geometry of Lemma \path{injective_union} is formalized
down to the selected complement and closure forms used by the two inverse
applications. The rectangular decompositions of $R$ and $S$ are formalized, and
the conditional derivations of injectivity for $R$ and $S$ from union closure
are also formalized. The two edge blocks removed from the displayed $T$-window
are injective under the coordinate rectangular hypotheses. The finite-lattice
complementary block around the chosen edge has a set-theoretic definition.
The exact-cover criterion by $2\times3$ and $3\times2$ rectangles is only a
restricted sufficient condition, and the origin-based instances above show
that it is not the source's injectivity argument. The source-faithful covers
and the every-edge blocking construction are not yet derived. The remaining
open gap is to prove the tensor-level union theorem, realign the local
$T$-injectivity argument with the finite PEPS geometry of the source, translate
the collar decomposition around each edge, and then prove the unconditional
injectivity of $R$, $S$, and the edge-complementary block.

The normal route itself is tracked by \ghissue{1373}. The normal-section
diagram work is tracked by \ghissue{1376}. The general PEPS status tracker is
\ghissue{1252}.

\section{Remaining mathematical obligations}

The normal PEPS theorem cannot be obtained by citing the injective theorem
alone. The lists below are the authoritative record for the square-lattice
normal blocking route. Lean docstrings and blueprint entries should point here
rather than repeat this status in full.

The following part of the formalization is already present.
\begin{itemize}
    \item Region injectivity is represented by
          \leanid{TNLean.PEPS.RegionInjectivityData}; the same infrastructure
          contains the region-complement definitions
          \leanid{TNLean.PEPS.regionComplement} and
          \leanid{TNLean.PEPS.regionOutsideUnion}, and the finite-region
          union-closure hypothesis
          \leanid{TNLean.PEPS.RegionInjectivityUnionClosure}.
    \item The four-region decomposition for Lemma \path{injective_union} is
          formalized both by named regions and by the indexed family
          \(A\setminus B\), \(A\cap B\), \(B\setminus A\), and
          \((A\cup B)^c\), including the reconstruction, disjointness, and
          selected-complement identities used by the two inverse applications.
    \item The three contraction maps for Lemma \path{injective_union} are
          separated as
          \leanid{TNLean.PEPS.InjectiveUnionContractionMaps}.
    \item The abstract contraction-chain criterion for Lemma
          \path{injective_union} is formalized by
          \leanid{TNLean.PEPS.InjectiveUnionContractionChain}: if the three
          implications
          \[
              \mathcal C_{\{0,1,2\}}(X)=0
              \Longrightarrow \mathcal C_{\{2\}}(X)=0
              \Longrightarrow \mathcal C_{\{1,2\}}(X)=0
              \Longrightarrow X=0
          \]
          are supplied, then the union contraction has trivial kernel.
    \item The square-lattice coordinate model, contiguous \(2\times3\) and
          \(3\times2\) rectangles, and the coordinate rectangular-injectivity
          hypotheses are present.
    \item The source regions \(R\) and \(S\) are explicit finite coordinate
          regions, their rectangular decompositions are formalized, and their
          injectivity follows conditionally from rectangular injectivity and
          union closure.
    \item The displayed local \(T\)-region, the finite-lattice
          edge-complement regions, their sufficient rectangular-cover criteria,
          the no-cover diagnostics for the present origin-based local models,
          and the normalized collar decompositions are formalized.
    \item The one-edge, translated-window, and margin-cover hypotheses
          for normal edge blocking are present, together with their endpoint,
          disjointness, coverage, and injective-chain consequences.
    \item The general normal PEPS blocking hypotheses
          \leanid{TNLean.PEPS.NormalPEPSBlockingHypotheses} provide their
          edge and one-site consequences directly, so downstream square-lattice
          arguments can use these hypotheses rather than restating the
          edge-blocking and one-site hypotheses separately.
\end{itemize}

The remaining obligations are the following.

\begin{enumerate}
    \item Prove the tensor-level union-of-injective-regions lemma. The proof
          must follow the source proof by applying the left inverses for $A$
          and $B$ in sequence. The disjoint case is now formalized over a bare
          three-block geometry (no red-block injectivity and no distinguished
          edge): for disjoint injective regions $S$ and $T$ with positive bond
          dimensions, $S\cup T$ is injective, and its nonempty finite iteration
          over pairwise-disjoint families is also available.\footnote{Formal
          declarations:
          \leanid{TNLean.PEPS.regionBlockedTensorInjective_union_disjoint} and
          \leanid{TNLean.PEPS.regionBlockedTensorInjective_biUnion_disjoint}.}
          The general overlapping statement
          $\operatorname{inj}(A)\wedge\operatorname{inj}(B)\Rightarrow
          \operatorname{inj}(A\cup B)$ remains open: the displayed regions $R$
          and $S$ are unions of \emph{overlapping} rectangles, so they require
          the four-region two-inverse descent rather than the disjoint engine.
    \item Realign the source proof of injectivity for the local and rotated
          \(T\)-regions with the finite PEPS geometry. The present exact-cover
          criterion is a useful sufficient condition, but the recorded
          point-forcing obstructions show that it is not the source argument
          for the origin-based local models.
    \item Prove the every-edge finite-geometry construction. The normalized and
          translated one-edge data give the desired red, blue, and
          complementary injective regions when the corresponding window or
          margin-cover datum is supplied; the remaining task is to construct
          source-faithful data for every edge, including boundary edges.
    \item Reuse the injective-chain insertion route after this blocking:
          \ghissue{1367}, \ghissue{1368}, \ghissue{1363}, \ghissue{1364},
          \ghissue{1361}, and \ghissue{1360}. The formal material for this route
          separates into unconditional region-insertion algebra, the
          identity-insertion bridge to the closed state, and a conditional
          matrix-transfer datum.

          The region-inserted coefficient is linear in the inserted matrix and,
          under the same positive-bond restriction as in the edge case, determines
          that matrix from all region and complement physical configurations.
          These facts use linear independence of the blocked weight family of
          \(R\) and of the complement \(V\setminus R\), and give the region
          analogue of the injectivity of the edge-inserted coefficient.\footnote{
          Formal declarations:
          \(\leanid{TNLean.PEPS.regionInsertedCoeff_add}\),
          \(\leanid{TNLean.PEPS.regionInsertedCoeff_smul}\),
          \(\leanid{TNLean.PEPS.regionInsertedCoeff_injective}\), and
          \(\leanid{TNLean.PEPS.edgeInsertedCoeff_injective}\).}
          Once an explicit edge matrix transfer has matched region-inserted
          coefficients, is multiplicative, and preserves the unit, it gives a
          region insertion-algebra isomorphism and hence, by Skolem--Noether, a
          per-edge gauge matrix.\footnote{
          Formal declarations:
          \(\leanid{TNLean.PEPS.RegionInsertionTransfer}\),
          \(\leanid{TNLean.PEPS.isRegionBlockedInsertionAlgebraIsomorphism_of_transfer}\), and
          \(\leanid{TNLean.PEPS.exists_regionEdgeGauge_of_transfer}\).}

          The identity-insertion bridge is the following equation. Here
          \(C_A(R,f,X;\sigma,\tau)\) denotes the region-inserted coefficient for
          the tensor \(A\), the region \(R\), the boundary edge \(f\), the boundary
          matrix \(X\), the region physical configuration \(\sigma\), and the
          complement physical configuration \(\tau\). The symbol \(D_A(e)\)
          denotes the bond dimension of the edge \(e\), \(\partial R\) denotes the
          set of boundary edges of \(R\), and
          \(\langle\sigma,\tau\mid\Psi_A\rangle\) denotes the closed-state
          coefficient of the assembled physical configuration:
          \[
              C_A(R,f,1;\sigma,\tau)
              =
              \Bigl(\prod_{e\notin \partial R} D_A(e)\Bigr)
              \langle \sigma,\tau \mid \Psi_A\rangle
          \tag{4.1}
          \]
          The proof decomposes the closed coefficient into region and complement
          contractions, rewrites identity insertion as a double sum over
          boundary-agreeing virtual configurations, and collapses that double sum
          by merging the two virtual configurations. The factor
          \(\prod_{e\notin\partial R}D_A(e)\) is therefore part of the formula,
          not a normalization that can be discarded.\footnote{
          Formal declarations:
          \(\leanid{TNLean.PEPS.prod_assembleRegionσ_split}\),
          \(\leanid{TNLean.PEPS.regionInsertedCoeff_identity_eq_doubleSum}\),
          \(\leanid{TNLean.PEPS.regionMerge}\),
          \(\leanid{TNLean.PEPS.regionProd_eq_merge}\),
          \(\leanid{TNLean.PEPS.complementProd_eq_merge}\),
          \(\leanid{TNLean.PEPS.regionInteriorBondProd}\),
          \(\leanid{TNLean.PEPS.stateCoeff_eq_regionComplement}\), and
          \(\leanid{TNLean.PEPS.regionInsertedCoeff_one_eq_stateCoeff}\).}
          Under matched bond dimensions the corresponding interior bond products
          for \(A\) and \(B\) agree, so the same factor is available in the
          matched-coefficient transfer.\footnote{
          Formal declarations:
          \(\leanid{TNLean.PEPS.regionInteriorBondProd_congr}\) and
          \(\leanid{TNLean.PEPS.regionInsertedCoeff_transfer_of_realizes}\).}

          For a second tensor \(B\), let \(T_f(X)\) denote the boundary matrix
          assigned to \(X\) on the edge \(f\) by the desired region insertion
          transfer. The formal candidate pulls the in-region endpoint operator of
          \(A\) back through \(B\)'s vertex and reads it off through a fixed
          residual reference frame, in analogy with the edge transfer matrix.
          The available recovery theorem is conditional: it reduces the matched
          coefficient identity to two hypotheses, namely that the transferred
          endpoint operator preserves the image of \(B\)'s local tensor map at the
          in-region vertex and that its virtual pullback through \(B\) has the
          single-edge insertion form on \(f\).\footnote{
          Formal declarations:
          \(\leanid{TNLean.PEPS.regionTransferMatrix}\),
          \(\leanid{TNLean.PEPS.edgeTransferMatrix}\), and
          \(\leanid{TNLean.PEPS.regionTransferMatrix_realizes_of_image}\).}
          These are the region analogues of the two consequences extracted by the
          edge physical-to-virtual recovery.\footnote{
          Formal declaration:
          \(\leanid{TNLean.PEPS.physical_to_virtual_insertion}\).}

          A useful image-preservation argument is available only in the
          vertex-injective regime. If \(B\) is vertex injective, then the
          state-vector coefficients at the in-region vertex span the local
          virtual coefficient space by row--column rank duality for the
          vertex-complement tensor family of \(B\). Under
          \(\operatorname{SameState}(A,B)\) this identifies the local tensor images
          of \(A\) and \(B\), and the transferred endpoint operator preserves that
          image.\footnote{
          Formal declarations:
          \(\leanid{TNLean.PEPS.span_stateOpenCoeff_eq_top}\),
          \(\leanid{TNLean.PEPS.vertexComplementTensorInjective_of_isVertexInjective}\),
          \(\leanid{TNLean.PEPS.range_localTensorMap_eq_span_stateVec}\),
          \(\leanid{TNLean.PEPS.range_localTensorMap_eq_of_sameState}\), and
          \(\leanid{TNLean.PEPS.regionInsertionOp_localProjectorAt_eq}\).}
          This does not discharge the normal-PEPS route of
          \cite[Section~3]{Molnar2018NormalPEPS}: the source theorem assumes
          injectivity after blocking, not single-vertex injectivity. The
          region-injective proof still requires the same spanning and
          image-preservation statement derived from blocked-region injectivity.

          Even in the vertex-injective setting, the full realization is
          conditional on the remaining matrix-structure transfer: the virtual
          pullback of the transferred endpoint operator must be the matrix
          insertion of \(T_f(X)^{\mathsf T}\) on \(f\).\footnote{
          Formal declaration:
          \(\leanid{TNLean.PEPS.regionTransferRealizesAt_of_hform}\).}
          Equivalently, the change of frame between \(A\)'s and \(B\)'s coefficient
          spaces at the in-region vertex must intertwine the single-bond
          insertions on \(f\). This is the region analogue of the incident-matrix
          form read from the edge resonate identity.

          This matrix-structure transfer is now reduced to a single structural
          fact. Because the region transfer matrix is, by definition, the read-off
          of the virtual pullback, the matrix-structure hypothesis holds exactly
          when that pullback is of single-bond insertion form on \(f\); the
          read-off then inverts the insertion and the round-trip closes the
          hypothesis. The image-preservation half is unconditional in the
          vertex-injective setting, so the pullback already realizes the
          transferred endpoint operator on every local tensor image of \(B\) at the
          in-region vertex. What remains is to show the pullback is of single-bond
          insertion form, which is the residual-independence the edge resonate
          identity supplies.\footnote{
          Formal declarations:
          \(\leanid{TNLean.PEPS.hform_of_isIncidentMatrixForm}\) and
          \(\leanid{TNLean.PEPS.localTensorMap_localVirtualOpOfPhysicalOpAt_regionInsertionOp}\).}

          The pin available from the closed-state transfer alone is non-circular but
          insufficient on its own. Evaluating the transferred endpoint operator of
          \(A\) on \(B\)'s state vector at the endpoint physical leg recovers
          \(A\)'s region-inserted coefficient; varying the
          region configuration at the endpoint vertex makes the leg range over the
          whole physical space while the state vector stays fixed, so the pin determines
          the endpoint operator's output on every state vector. But this only equates
          the operator with the read-off matrix \emph{at the single reference residual
          configuration} used to define \(T_f\); extending it to all residual
          configurations is the residual-independence the read-off must satisfy, and
          a single pin per region configuration cannot supply it. The closed-state
          transfer reads only one bond contraction (through \(R\)'s endpoint vertex),
          so it cannot reconcile the two residual frames the way the edge proof does.\footnote{
          Formal declarations:
          \(\leanid{TNLean.PEPS.regionInsertionOp_regionStateVec_pin}\) and
          \(\leanid{TNLean.PEPS.regionStateVec_update_vmem}\).}

          The faithful route mirrors the edge proof's two-endpoint structure with
          \(R\) and its set complement \(\mathrm{univ}\setminus R\) as the two contracted
          blocks of the single boundary edge \(f\): the in-region endpoint vertex
          plays the role of one edge endpoint, the out-of-region endpoint vertex of
          \(f\) (the in-region endpoint of \(f\) for the complement region) plays the
          other, and the empty interior of the single edge \(f\) makes the middle
          block trivial. The region-inserted coefficient already carries this doubled
          boundary-configuration sum (one configuration read by \(R\), one by the
          complement, coupled by the inserted matrix on \(f\)). The missing steps are:
          (i) a complement-side realization, the analogue of
          the realization of the region-inserted coefficient applied to the complement
          region with \(f\) as its boundary edge; (ii) a region resonate identity
          equating the \(v\)-side and complement-side
          realizations of \(A\)'s region-inserted coefficient, transferred to \(B\) by
          \(\operatorname{SameState}\); (iii) inverting both endpoints through the
          blocked-region left inverses (using the blocked-region
          injectivity of \(R\) and of the complement); and (iv) a region reconcile
          forcing
          the two read-off matrices to coincide. This is the region-granularity port
          of the edge insertion-realization argument and remains the open
          obligation. Once it is supplied, the transfer datum and the per-edge gauge
          follow mechanically.

          Step (i) now has a cast-free reading. The complement-side reading is the
          region-inserted coefficient read with the region and its set complement
          exchanged and the inserted matrix transposed on \(f\). The exchange is the
          block-swap symmetry of the abstract two-block inserted coefficient, which
          stays over \(R\)'s boundary-edge index set and so avoids the dependent-type
          reindexing \(\mathrm{univ}\setminus(\mathrm{univ}\setminus R)=R\). What still
          requires that reindexing is factoring the complement block through the
          out-of-region endpoint vertex's local tensor map for the inversion of step
          (iii).\footnote{
          Formal declarations:
          \(\leanid{TNLean.PEPS.twoBlockInsertedCoeff_swap}\) and
          \(\leanid{TNLean.PEPS.regionInsertedCoeff_eq_complementTwoBlock}\).}

          Steps (ii) and the inversion infrastructure of step (iii) are now
          formalized. The out-of-region-endpoint pin reads the region-inserted
          coefficient through the out-of-region endpoint vertex by instancing the
          in-region pin at the set complement and casting back through the
          complement-side cast identity. Equating it with the in-region-endpoint pin
          gives the region resonate identity: the in-region and out-of-region endpoint
          operators of the first tensor, applied to the second tensor's closed state
          vectors, agree at the two endpoint legs. The region-inserted coefficient
          factors, as a function of the complement (resp. region) physical
          configuration, through the blocked-region tensor map of the set complement
          (resp. of the region); each blocked-region left inverse reads off the
          corresponding row function. What remains is the reconcile of step (iv):
          exhibiting the first tensor's region-inserted coefficient, as a function of
          the region physical configuration, in the image of the second tensor's
          region blocked tensor map (the membership that the resonate transfer across
          \(\operatorname{SameState}\) supplies), then matching the two read-off row
          functions and concluding incident-matrix form of the virtual pullback
          through \(\leanid{TNLean.PEPS.hform_of_isIncidentMatrixForm}\).

          The remaining content of step (iv) is now isolated to a single predicate.
          Incident-matrix form of the virtual pullback, for every inserted matrix,
          assembles the realization bundle and hence the region insertion transfer
          datum and the per-edge gauge, all unconditionally on top of that predicate.
          The predicate names exactly the reconcile: the virtual pullback of the
          transferred in-region endpoint operator is of single-bond insertion form on
          the boundary edge. With it in both directions the per-edge gauge follows
          mechanically, so the open obligation is reduced to establishing this
          predicate from the region resonate identity and the two blocked-endpoint
          inversions.

          The reconcile predicate is in turn reduced to a single coefficient
          identity. The leg-wise pin reads the in-region endpoint operator of the
          first tensor, applied to the second tensor's closed state vector, leg by
          leg as the first tensor's region-inserted coefficient scanned over the
          endpoint physical configuration. The closed state-vector coefficients span
          the full local virtual coefficient space at the in-region endpoint vertex,
          so a realization on the state vectors extends to all coefficients. Together
          these show: if for every inserted matrix on the first tensor's bond there
          is a matrix on the second tensor's bond whose region-inserted coefficient
          equals the first's at every physical configuration, then the reconcile
          predicate holds, with the read-off matrix the transpose of the matching
          matrix. The whole remaining obligation is therefore exactly this
          coefficient transfer: producing, for each inserted matrix, a matching
          matrix on the second tensor. This is the content the region resonate
          identity and the two blocked-endpoint inversions must supply; it is
          equivalent to the insertion-algebra isomorphism of the source lemma.\footnote{
          Formal declarations:
          \(\leanid{TNLean.PEPS.regionInsertionOp_regionStateVec_pin_leg}\),
          \(\leanid{TNLean.PEPS.regionInsertionOp_realizes_of_realizes_on_stateVec}\),
          \(\leanid{TNLean.PEPS.regionInsertionOp_regionStateVec_eq_of_coeff_eq}\),
          \(\leanid{TNLean.PEPS.isIncidentMatrixForm_of_coeff_eq}\), and
          \(\leanid{TNLean.PEPS.regionResonateReconcile_of_coeff_transfer}\).}

          The coefficient transfer is itself reduced to a single structural
          statement about a doubly-indexed transfer coefficient. The first tensor's
          region-inserted coefficient factors through both blocked endpoints of the
          second tensor at once: the v-side complement factorization gives the
          coefficient as a function of the complement physical configuration, and
          the symmetric region factorization (read off the complement-side cast
          identity and transported across the double complement) gives it as a
          function of the region physical configuration. Combining the two, the
          coefficient is the double sum of a transfer coefficient against the second
          tensor's region and complement blocked weights, with the transfer
          coefficient read off by the two blocked-endpoint left inverses; the
          required membership in the second tensor's complement block image holds
          because the region row is a finite linear combination of the v-side
          complement-image coefficients. Once the transfer coefficient is shown to
          have the incident-matrix coupling form of a single matrix on the boundary
          bond --- coupling only the boundary-edge legs and contracting the residual
          legs by the identity --- that matrix is the matching matrix the
          coefficient transfer asks for. This incident-matrix form of the transfer
          coefficient is the only remaining content; it is the region analogue of
          the source step \(V=W\), forced by the resonate identity together with the
          full injectivity of both blocked endpoints.

          The incident-matrix form of the transfer coefficient is now reduced, by a
          sorry-free chain, to the incident-matrix form of the virtual pullback
          itself. The transfer-coefficient column at a fixed complement boundary
          configuration is the second tensor's region blocked left inverse of the
          v-side row. The v-side row is the first tensor's in-region endpoint
          operator, from the inserted matrix, applied to the second tensor's region
          weight vector; since that weight vector is the second tensor's local tensor
          map of the region open coefficient, the unconditional image-preservation
          realization rewrites the v-side row as the second tensor's local tensor map
          of the virtual pullback applied to the region open coefficient. If the
          virtual pullback is of incident-matrix form, the second tensor's own
          in-region endpoint operator of that matrix realizes it, and the inner-sum
          realization expands the v-side row to the incident-matrix sum against the
          region blocked weights; the region blocked left inverse then collapses each
          weight to the standard basis configuration, reading off the incident-matrix
          coupling of the transfer coefficient. The reconcile predicate
          \(\leanid{TNLean.PEPS.RegionResonateReconcile}\) is exactly this
          incident-matrix form of the virtual pullback for every inserted matrix, so
          the whole open obligation is now this single predicate. Establishing it
          requires the out-of-region (complement-vertex) inversion of the region
          resonate identity, the region analogue of the right-endpoint inversion that
          supplies the residual-independence of the read-off matrix at the edge
          level; the in-region pin alone fixes the matrix only at the reference
          residual.\footnote{
          Formal declarations (the sorry-free reduction):
          \(\leanid{TNLean.PEPS.transferCoeff_column_eq_regionBlockedLeftInverse_vSideRow}\),
          \(\leanid{TNLean.PEPS.vSideRow_eq_localTensorMap_virtualPullback}\),
          \(\leanid{TNLean.PEPS.vSideRow_eq_incidentSum_of_virtualPullback_incidentForm}\),
          and
          \(\leanid{TNLean.PEPS.transferCoeff_eq_incidentForm_of_virtualPullback_incidentForm}\).}\footnote{
          Formal declarations:
          \(\leanid{TNLean.PEPS.regionInsertedCoeff_eq_complement_blockedMap_B}\),
          \(\leanid{TNLean.PEPS.regionInsertedCoeff_eq_region_blockedMap_B}\),
          \(\leanid{TNLean.PEPS.regionBlockedTensorMap_doubleCompl}\),
          \(\leanid{TNLean.PEPS.regionRowB_mem_range}\),
          \(\leanid{TNLean.PEPS.regionInsertedCoeff_eq_doubleSum_transferCoeff}\),
          \(\leanid{TNLean.PEPS.vSideRow_eq_region_blockedMap_transferCoeff}\), and
          \(\leanid{TNLean.PEPS.regionInsertedCoeff_eq_of_transferCoeff_form}\).}\footnote{
          Formal declarations:
          \(\leanid{TNLean.PEPS.RegionResonateReconcile}\),
          \(\leanid{TNLean.PEPS.regionTransferRealizes_of_isIncidentMatrixForm}\),
          \(\leanid{TNLean.PEPS.regionTransferRealizes_of_reconcile}\),
          \(\leanid{TNLean.PEPS.regionInsertionTransfer_of_reconcile}\), and
          \(\leanid{TNLean.PEPS.exists_regionEdgeGauge_of_reconcile}\).}\footnote{
          Formal declarations:
          \(\leanid{TNLean.PEPS.regionInsertionOp_regionStateVec_pin_compl}\),
          \(\leanid{TNLean.PEPS.region_resonate_identity}\),
          \(\leanid{TNLean.PEPS.regionInsertedCoeff_eq_complement_blockedMap}\),
          \(\leanid{TNLean.PEPS.regionInsertedCoeff_eq_region_blockedMap}\),
          \(\leanid{TNLean.PEPS.regionBlockedLeftInverse_complement_regionInsertedCoeff}\), and
          \(\leanid{TNLean.PEPS.regionBlockedLeftInverse_region_regionInsertedCoeff}\).}\footnote{
          Formal declarations:
          \(\leanid{TNLean.PEPS.regionInsertedCoeff}\),
          \(\leanid{TNLean.PEPS.regionInsertedCoeff_eq_regionRealizationSum}\),
          \(\leanid{TNLean.PEPS.resonate_invert_right_endpoint}\),
          \(\leanid{TNLean.PEPS.resonate_invert_left_endpoint}\),
          \(\leanid{TNLean.PEPS.resonate_endpoint_coeff_reconcile}\),
          \(\leanid{TNLean.PEPS.InsertionRealization}\),
          \(\leanid{TNLean.PEPS.regionInsertionTransfer_of_realizes}\), and
          \(\leanid{TNLean.PEPS.exists_regionEdgeGauge_of_transfer}\).}

          The reconcile above is also reachable in a \emph{block frame} that uses no
          single-vertex injectivity. The foundational input is the block-level image
          coincidence: under \(\operatorname{SameState}\), with the complement block
          blocked-tensor injective and positive bond dimensions, the range of the
          region blocked tensor map is the same for the two tensors, because both
          ranges equal the span of the partial states across the region cut, and the
          partial states depend only on the closed state coefficient. The same fact
          for the complement block follows by transporting blocked-tensor injectivity
          across the double complement. With these two image coincidences, the first
          tensor's region-inserted coefficient, read as a function of the region
          physical configuration, lies in the second tensor's region block image
          (region side), and the transferred row, as a function of the complement
          physical configuration, lies in the second tensor's complement block image
          (complement side); the two blocked left inverses read off the same transfer
          coefficient as in the vertex frame, and the first tensor's coefficient is
          the double sum of that transfer coefficient against the second tensor's two
          blocked weights. None of this uses single-vertex injectivity.

          In the block frame the remaining content is isolated cleanly. By double
          blocked injectivity of the second tensor, the double-sum representation of a
          coefficient by a transfer coefficient is unique, so the transfer
          coefficient of an inserted matrix has the incident-matrix coupling form of a
          single bond matrix if and only if the first tensor's region-inserted
          coefficient of that inserted matrix equals the second tensor's of that bond
          matrix. The whole open obligation is therefore equivalent to the coefficient
          transfer alone --- for every inserted matrix on the first tensor's bond, a
          matching bond matrix on the second tensor --- with the matching matrix
          carried directly by the equivalence. The identity insertion is matched
          across \(\operatorname{SameState}\), so the transfer coefficient of the
          identity is already the incident-matrix coupling of the identity. The
          residual open content is the incident-matrix coupling form of the transfer
          coefficient for a general inserted matrix: the transfer coefficient must be
          local on the boundary bond, depending on the two boundary configurations
          only through their bond legs. This locality is the bond-level content of the
          source step \(V=W\) and is the one fact the block-frame foundation does not
          supply, because the block image coincidence preserves only the full weight
          span and not the per-leg decomposition across the boundary bond.\footnote{
          Formal declarations (block frame, sorry-free):
          \(\leanid{TNLean.PEPS.range_regionBlockedTensorMap_eq_of_sameState}\),
          \(\leanid{TNLean.PEPS.regionBlockedTensorInjective_doubleCompl}\),
          \(\leanid{TNLean.PEPS.range_regionBlockedTensorMap_compl_eq_of_sameState}\),
          \(\leanid{TNLean.PEPS.regionInsertedCoeff_eq_blockTransferRow}\),
          \(\leanid{TNLean.PEPS.blockTransferRow_mem_range}\),
          \(\leanid{TNLean.PEPS.regionInsertedCoeff_eq_doubleSum_transferCoeff_block}\),
          \(\leanid{TNLean.PEPS.doubleSum_kernel_injective}\),
          \(\leanid{TNLean.PEPS.transferCoeff_eq_incidentKernel_iff_coeff_eq}\),
          \(\leanid{TNLean.PEPS.transferCoeff_one_eq_incidentKernel_one}\),
          \(\leanid{TNLean.PEPS.coeffTransfer_iff_transferCoeff_incidentForm}\), and
          \(\leanid{TNLean.PEPS.exists_regionEdgeGauge_of_blockCoeffTransfer}\).}

          The block-frame foundation is now assembled around the edge-blocking
          triple. The union of two injective blocks is injective; applied to the blue
          and complement blocks of a one-edge blocking datum, this gives injectivity of
          the host block, the set complement of the red block, which every block-frame
          backbone lemma needs at the complement of the red region. With it, the four
          block injectivities of the two tensors at the red region and at its host are
          all read off a shared one-edge blocking datum, recording that the two tensors
          are blocked around the same edge into the same red, blue, and complement
          regions, together with positive bond dimensions. On top of these, the
          per-edge gauge is assembled from one predicate alone: bond locality of the
          transfer coefficient on the boundary edge, in both directions. Bond locality
          is proved equivalent to the coefficient transfer, so the open content of the
          block frame is now exactly this single predicate, with the per-edge gauge and
          all four injectivities derived from it and the shared blocking datum without
          single-vertex injectivity. The identity transfer of a single tensor is
          bond-local, the diagonal anchor of the cross-tensor transfer. The remaining
          open obligation is bond locality of the cross-tensor transfer for a general
          inserted matrix, the block-level content of the step \(V=W\); it is the only
          fact between the assembled foundation and the per-edge gauge.\footnote{
          Formal declarations (block frame, sorry-free):
          \(\leanid{TNLean.PEPS.regionBlockedTensorInjective_union}\),
          \(\leanid{TNLean.PEPS.regionBlockedTensorInjective_compl_red}\),
          \(\leanid{TNLean.PEPS.SharedNormalEdgeBlockingData}\),
          \(\leanid{TNLean.PEPS.IsBondLocalTransferKernel}\),
          \(\leanid{TNLean.PEPS.bondLocal_iff_coeffTransfer}\),
          \(\leanid{TNLean.PEPS.coeffTransfer_of_bondLocal}\),
          \(\leanid{TNLean.PEPS.isBondLocalTransferKernel_self}\), and
          \(\leanid{TNLean.PEPS.SharedNormalEdgeBlockingData.exists_regionEdgeGauge}\).}

          The block-frame transfer kernel \(K_M(\mu,\nu')\) is the double blocked
          left inverse of the first tensor's region-inserted coefficient of \(M\)
          through the second tensor's region and complement blocks. Its bond locality
          is the incident-matrix coupling form \(K_M=\text{incidentKernel}(N)\): the
          kernel couples the two boundary configurations only through their legs on
          the boundary edge \(f\). The criterion is now stated through the \(f\)-leg
          split of a boundary configuration: \(K_M\) is the incident kernel of \(N\)
          exactly when, read through the split, it vanishes off the diagonal of the
          two residual boundary configurations and reads \(N\) on the matching \(f\)-leg.
          The transferred row that feeds the kernel is the second tensor's region
          blocked left inverse of the first tensor's coefficient, which is exactly the
          \(A\leftrightarrow B\) region basis change applied to the first tensor's own
          region row \(\text{rowInsert}_f(M)\) of the complement weight row --- a row
          that is already bond-\(f\) local. Bond locality of \(K_M\) is therefore
          equivalent to the region basis change preserving the bond-\(f\)/residual
          decomposition, that is, conjugating the first tensor's \(f\)-leg row
          insertion of \(M\) to the second tensor's \(f\)-leg row insertion of a single
          \(N\). This is the basis-change intertwining the cross-tensor reduction
          consumes.\footnote{
          Formal declarations (block frame, sorry-free):
          \(\leanid{TNLean.PEPS.incidentKernel_eq_split}\),
          \(\leanid{TNLean.PEPS.transferCoeff_eq_incidentKernel_iff_split}\),
          \(\leanid{TNLean.PEPS.blockTransferRow_eq_basisChange_regionRegionRow}\),
          \(\leanid{TNLean.PEPS.regionInsertedCoeff_eq_of_basisChange_intertwine}\),
          \(\leanid{TNLean.PEPS.basisChange_intertwine_iff_coeffTransfer}\), and
          \(\leanid{TNLean.PEPS.isBondLocalTransferKernel_of_coeffTransfer}\).}

          The port of the edge resonate inversion to this block granularity hits one
          structural mismatch that is not a matter of length. The edge inversion
          \(\leanid{TNLean.PEPS.resonate_invert_right_endpoint}\) consumes the
          bond-contracted resonate identity
          \(\leanid{TNLean.PEPS.resonate_middle_inverted}\), which pairs the
          left-endpoint and right-endpoint physical operators \emph{indexed by the
          shared bond} \(k\): the two endpoints are two single vertices sharing one
          bond, and the two operators act on the same physical space \(\mathbb C^d\).
          In the block frame the two endpoints are the red region and its complement
          sharing the boundary edge \(f\), and the only available factorizations of the
          first tensor's coefficient through the second tensor's complement block in the
          bond-\(f\)-paired form needed by the inversion --- the v-side and
          \(\sigma\)-side blocked maps
          (\(\leanid{TNLean.PEPS.regionInsertedCoeff_eq_complement_blockedMap_B}\),
          \(\leanid{TNLean.PEPS.regionInsertedCoeff_eq_region_blockedMap_B}\)) ---
          carry the single-vertex hypothesis \(\text{LinearIndependent}\,(A_v)\) at the
          endpoint vertex of \(f\). The block-image coincidence
          \(\leanid{TNLean.PEPS.range_regionBlockedTensorMap_eq_of_sameState}\) supplies
          the full-weight-span identity through the second tensor's whole region block
          as a single basis change, but it does not supply the per-\(f\)-leg bond
          pairing the inversion consumes; the only block-only complement-side reading
          of the coefficient is through the double left inverse \(K_M\) itself, which is
          the very kernel whose bond locality is the goal. The block-frame port
          therefore requires a vertex-injectivity-free bond-\(f\)-paired complement
          factorization of the coefficient before the edge inversion structure
          transfers; constructing it from the two block realizations through the second
          tensor's partial states
          (\(\leanid{TNLean.PEPS.regionInsertedCoeff_eq_blockRealizeOp_regionPartialState_B}\),
          \(\leanid{TNLean.PEPS.regionInsertedCoeff_eq_blockRealizeOp_complementPartialState_B}\))
          is the remaining open step.
    \item In the translationally invariant case, show that the edge gauges obtained
          after blocking reduce to one horizontal gauge $X$ and one vertical gauge
          $Y$, up to scalar.
    \item Prove the scalar condition $\lambda^{nm}=1$ for the square-lattice
          theorem, and state separately the scalar freedom in the non-translation
          invariant normal theorem.
\end{enumerate}

\section{Source-to-formalization ledger}

\begin{itemize}
    \item Lemma \path{injective_union}: blueprint
          \path{lem:peps_injectiveRegion_union}, issue \#1374.
    \item Finite iteration of Lemma \path{injective_union}: blueprint
          \path{lem:peps_regionInjectivityUnionClosure_finite}, issue \#1374.
    \item Regions $R$, $S$, and $T$: blueprint
          \path{def:peps_normalSquareBlockingRegions}, issue \#1375. This
          node also records the conditional implication from coordinate
          rectangular injectivity, union closure, and a rectangular cover of
          the displayed $T$-region to the abstract $R,S,T$ blocking-region
          structure.\footnote{The formal declaration is
          \leanid{TNLean.PEPS.normalSquareBlockingRegions_of_TCover}.}
    \item Injectivity of the two edge blocks removed from the displayed
          $T$-window, and of their union under union closure: blueprint
          \path{lem:peps_normalSquareRegionT_edgeBlocks_injective}, issue
          \#2004.
    \item General rectangular-cover criterion for square-lattice regions:
          blueprint \path{lem:peps_normalSquareRectangleCover_injective}, issue
          \#2004.
    \item Rectangular-cover criterion for injectivity of the displayed
          $T$-region: blueprint
          \path{lem:peps_normalSquareRegionT_injective_of_cover}, issue
          \#2004. This node also contains the rotated vertical local
          $T$-region cover criterion.
    \item Obstruction to applying that criterion directly to the present
          origin-based local-window $T$ models and edge-complement models:
          blueprint
          \path{lem:peps_normalSquareRegionT_no_five_by_six_cover}, issue
          \#2004.
    \item Finite-lattice complementary block around an edge: blueprint
          \path{lem:peps_normalSquareEdgeComplementRegion}, issue \#2004.
    \item Rectangular-cover criterion for injectivity of that block: blueprint
          \path{lem:peps_normalSquareEdgeComplementRegion_injective_of_cover},
          using the cover definition
          \path{def:peps_normalSquareRectangleCover}, issue \#2004.
    \item Extension of a local $T$ rectangular cover to a rectangular cover of
          the normalized horizontal-edge complementary block by adjoining the
          two top-collar rectangles: blueprint
          \path{lem:peps_normalSquareEdgeComplementRegion_coverOfT}, issue
          \#2004.
    \item Normalized $5\times7$ top-collar decomposition of that block:
          blueprint
          \path{lem:peps_normalSquareEdgeComplementRegion_topCollar}, issue
          \#2004. This node also contains a normalized $7\times5$
          collar identity for the analogous rectangular frame and the rotated
          vertical edge-complement collar identity.
    \item Normalized horizontal and vertical red/blue/complement blocking data:
          blueprint
          \path{lem:peps_normalSquareHorizontalEdge_blockingData}, issue
          \#2004. These use the common one-edge blocking-data structure.
          They also expose the corresponding injective-chain conclusions. The
          horizontal complement injectivity is reduced to a local $T$-cover,
          and the vertical complement injectivity is reduced to a rotated
          local $T$-cover.
    \item Translated horizontal red/blue/complement blocking data:
          blueprint
          \path{lem:peps_normalSquareHorizontalTranslatedEdge_blockingData},
          issue \#2004. This is the coordinate-origin-parametric horizontal
          step toward the every-edge construction. Its coordinate origin
          specializes to the normalized horizontal picture.
    \item Translated vertical red/blue/complement blocking data:
          blueprint
          \path{lem:peps_normalSquareVerticalTranslatedEdge_blockingData},
          issue \#2004. This is the coordinate-origin-parametric vertical
          step toward the every-edge construction. Its coordinate origin
          specializes to the normalized vertical picture.
    \item Conditional construction from translated edge windows:
          blueprint
          \path{thm:peps_normalSquareTranslatedEdgeBlockingHypotheses_of_windows},
          issue \#2004. This records the formal boundary before proving that
          every edge in the finite square lattice admits such a translated
          window. Coordinate right and upward edges are tied to the
          corresponding translated windows, and the margin hypotheses for
          actual coordinate right/up edges are explicit. Arbitrary horizontal
          and vertical edges inherit these criteria from their coordinate
          representatives. The oriented margin-and-cover data structure records
          the remaining per-edge input in the open rectangular coordinate
          model. One such datum directly supplies one-edge blocking data
          together with the corresponding endpoint, injectivity, disjointness,
          and coverage conclusions, and a family of such data determines the
          normal edge-blocking hypotheses. The resulting hypotheses recover the
          supplied datum and the same combined edge-blocking conclusions at
          each edge. The horizontal and vertical margin
          predicates are named, the coordinate right/up and oriented-edge
          window constructors accept those predicates directly, translated
          windows on coordinate right/up edges force those predicates, and the
          open $7\times7$ graph does not admit a family of these
          translated windows over all edges; four explicit boundary examples
          are recorded as a single conjunction. The per-edge margin-cover structure
          stores the same predicates, and the resulting open-boundary obstruction for
          horizontal edges without enough left or right margin and vertical
          edges without enough bottom or top margin is also formalized,
          including explicit $7\times7$ examples and their four-sided
          conjunction. Thus the present margin criterion is a sufficient interior
          criterion rather than the final every-edge construction.
          Consequently, the nonexistence of a family of this margin-cover data
          over all edges of the current open $7\times7$ graph is also
          formalized.
    \item A family of one-edge blocking data determining the uniform
          every-edge hypotheses, including the combined endpoint, injectivity,
          disjointness, and coverage assertion: blueprint
          \path{thm:peps_normalEdgeBlockingHypotheses_injectiveChain}, issue
          \#2004.
    \item Edge blocking into red, blue, and complementary injective regions:
          blueprint \path{thm:peps_normalEdgeBlockingToInjectiveChain}, issue \#1375.
    \item Application of Lemma \path{inj_isomorph} after normal blocking:
          downstream issues \#1367, \#1368, \#1363, and \#1364.
    \item Translationally invariant gauge absorption:
          blueprint \path{thm:peps_tiNormalGaugeAbsorption}, issue \#1375.
    \item Final $R$ and $S$ comparison:
          blueprint \path{thm:peps_oneVertexComplementComparison}, issue \#1360.
    \item Theorem \path{3}: blueprint
          \path{thm:fundamentalTheorem_normalSquarePEPS}, issue \#1375.
    \item General theorem \path{normal}: blueprint
          \path{thm:fundamentalTheorem_normalPEPS}, issue \#1375.
\end{itemize}

\section{Blueprint diagrams}

The normal route has its own tensor-network diagrams. These are not substitutes
for the injective theorem diagrams; they attach the additional normal blocking
layer to named proof nodes.\footnote{Chapter~13a calls the public TikZ commands
    in \path{blueprint/src/macros/tn_print.tex}. The web blueprint renders
    cached SVGs from the same commands through
    \path{blueprint/src/Packages/tn_diagrams.py}.}

\begin{itemize}
    \item Lemma \path{injective_union} shows the four regions
          $A\setminus B$, $A\cap B$, $B\setminus A$, and $(A\cup B)^c$.
    \item Definition \path{peps_normalSquareBlockingRegions} shows the regions
          $R$, $S$, and $T$ from the square-lattice proof.
    \item Lemma \path{peps_normalSquareEdgeComplementRegion_topCollar} shows the
          finite-lattice complementary block $A_3$ as the local $T$-region
          together with the two top-collar rectangles in the normalized
          $5\times7$ frame.
    \item Theorem \path{peps_normalEdgeBlockingToInjectiveChain} shows the
          red/blue/complementary blocking around a distinguished edge and the
          resulting three-site chain.
    \item Theorem \path{peps_tiNormalGaugeAbsorption} shows the final
          translationally invariant local gauge formula with horizontal gauge
          $X$ and vertical gauge $Y$.
\end{itemize}

\section{Conclusion}

The normal PEPS route is separated from the injective PEPS route in the
blueprint and in the GitHub issue tree. The formal proof still remains open:
the region-injectivity vocabulary, the four-region decomposition, and the
abstract contraction-chain criterion are present, but the tensor-level proof of
Lemma \path{injective_union} is not. The next genuine mathematical steps are
that tensor-level union lemma, the source-faithful local and rotated
$T$-injectivity argument, and the every-edge finite-geometry construction.
After these normal blocking hypotheses are proved, the downstream proof should
reuse the injective edge-insertion route rather than duplicate it. The
translationally invariant gauge absorption and the scalar condition remain the
final steps of the normal theorem.

\section{The vertex-frame obstruction and the block-frame route}

The region insertion machinery formalized so far inserts the comparison matrix
at the \emph{in-region endpoint vertex} $v$ of the boundary edge $f$ and
realizes it through $v$'s tensor (\path{regionInsertionOp}). With this frame the
coefficient transfer reduces cleanly to a single residual-frame statement: the
$v$-side row of the first tensor is of single-bond insertion form
(\path{coeffTransfer_of_vSideRow_incidentForm}). Every step of that reduction uses
only blocked-region injectivity of the region and its complement
(\path{regionInsertedCoeff_eq_of_vSideRow_incidentForm},
\path{transferCoeff_eq_incidentForm_of_vSideRow_incidentForm},
\path{doubleBlockedResonate}, \path{complSideRow_eq_complement_blockedMap_transferCoeff}).

The remaining $v$-side row statement, however, does not close under blocked-region
injectivity alone. Discharging it evaluates the first tensor's endpoint operator,
which factors as
$\Phi^A_v \circ (\text{incident insertion}) \circ \Psi^A_v$ with $\Psi^A_v$ a left
inverse of $\Phi^A_v=\path{localTensorMap}\ A\ v$, on the \emph{second} tensor's
region weight vectors. This requires the cross-tensor endpoint image inclusion
$\operatorname{range}(\Phi^B_v)\subseteq\operatorname{range}(\Phi^A_v)$, i.e.\ the
two single-vertex tensor images at $v$ coincide. That coincidence is available only
from single-vertex spanning (\path{span_stateOpenCoeff_eq_top}, derived from
\path{vertexComplementTensorInjective_of_isVertexInjective}), which is genuine
\emph{vertex} injectivity. Blocked-region injectivity is linear independence of the
blocked weight family indexed by boundary configurations; for a region with more
than one site it does not force the single-vertex marginal at $v$ to be injective,
and two tensors with the same global state can have distinct $v$-column spaces. The
disanalogy with the edge argument is structural: the edge resonate identity is
single-tensor (the abstract operators act on one family's components), whereas the
region resonate identity reads the first tensor's operator on the second tensor's
vectors, introducing a cross-tensor image mismatch absent at the edge level.

Hence the vertex-endpoint frame closes only the \emph{vertex-injective}
specialization, where the injective Fundamental Theorem already supplies the gauge
(\path{fundamentalTheorem_normalPEPS_of_vertexInjective}). The source's actual
argument is block-framed: the injective blocks play the role of super-sites, and the
insertion is realized on the \emph{block's} open boundary leg through the blocked
tensor, so the image-preservation step becomes the block-level range coincidence
$\operatorname{range}(\text{blocked map of }A)=\operatorname{range}(\text{blocked map of }B)$,
which is provable from blocked-region injectivity together with the same-state
hypothesis (the blocked-weight analogue of
\path{range_localTensorMap_eq_of_sameState}). The conditional chain from the
coefficient transfer to the per-edge gauge is already block-injective-clean
(\path{exists_regionEdgeGauge_of_coeffTransfer}, no single-vertex injectivity), so a
block-framed proof of the coefficient transfer feeds it directly. The block-framed
realization---a block-leg insertion operator with its own block-level resonate
inversion---is the genuine remaining obligation for the general region-injective
theorem and the recommended next development.

\bibliographystyle{alpha}
\bibliography{references}

\end{document}
